% ============================================================
\section{Experiments}
% ============================================================



\subsection{VLM and Computer-Using Agent Baselines}
We conduct human evaluations with annotators who possess basic software usage skills but have not been exposed to the tasks before. For each human trajectory, we record the number of interaction steps and the wall-clock completion time, and aggregate the measurements by platform. The human experiment details are described in Appendix~\ref{app:human_experiment}.
Then, we conduct experiments on (a) generalize representative models GPT-5.4, GPT-5.5~\citep{openai2026gpt54, openai2026gpt55}, Gemini-3.1-Pro, Gemini-3-Pro~\citep{googledeepmind2026gemini31pro,googledeepmind2025gemini3pro}, Claude Opus 4.7~\citep{anthropic2026opus47}, Qwen-3.6-Plus~\citep{alibabacloud2026qwen36plus}, and Kimi-K2.6~\citep{moonshot2026kimi26}, (b) specialized computer-using agents OpenCUA-72B~\citep{wang2025opencua}, ScaleCUA-32B~\citep{liu2025scalecua}. By default, agents receive the current screenshot and the three most recent interaction steps and screenshots. We request agents to complete the tasks within a max step limit of 100, and per-step time limit of 120 seconds. The prompts used for each platform in the experiments are provided in Appendix~\ref{app:prompting_details}. We present a summary of the results in Table~\ref{tab:main_results}.


\noindent
\textbf{The \textit{testmini} Subset.} \dataset comprises 1,200 high-quality real-world CUA tasks. To streamline evaluation and accelerate agent development validation, we extract a smaller representative subset named \textit{testmini} including 300 problems. The construction of \textit{testmini} involved a proportional random sampling strategy across different platforms. The quantitative evaluations in all subsequent experiments were conducted on this \textit{testmini} subset.
%, while the full set of results on all 1,000 problems are provided in Appendix~\ref{}.



\begin{table*}[t]
\centering
\scriptsize
\setlength{\tabcolsep}{3pt}
\renewcommand{\arraystretch}{1.18}
\definecolor{modelgroup}{HTML}{F0F0F0}
\definecolor{bestavg}{HTML}{FEE4E2}
\definecolor{secondavg}{HTML}{D1E9FF}
\definecolor{badgeClaude}{HTML}{D97706}
\definecolor{badgeScale}{HTML}{4F46E5}
\definecolor{badgeOpenCUA}{HTML}{0EA5E9}
\newcommand{\osmodelicon}[2][1.05em]{%
  \IfFileExists{figures/icons/#2}{%
    \raisebox{-0.18\height}{\includegraphics[height=#1]{figures/icons/#2}}\,%
  }{}%
}
\newcommand{\modelbadge}[2]{%
  \begingroup
  \setlength{\fboxsep}{1.1pt}%
  \raisebox{0.04em}{\colorbox{#1}{\textcolor{white}{\scriptsize\bfseries\strut #2}}}%
  \endgroup\,%
}
\caption{\textbf{Main results across \dataset platforms.} Each cell reports task success rate (\%) by platform and difficulty. \textbf{Avg.} is computed over available platform--difficulty cells for each model.}
\label{tab:main_results}
\resizebox{\textwidth}{!}{%
\begin{tabular}{>{\raggedright\arraybackslash}p{3.1cm}*{13}{c}}
\toprule
\multirow{3}{*}{\textbf{Model / Agent}}
& \multicolumn{13}{c}{\textbf{Success Rate ($\uparrow$)}} \\
\cmidrule(lr){2-14}
& \multicolumn{2}{c}{\windowsicon\,\textbf{Windows}}
& \multicolumn{2}{c}{\linuxicon\,\textbf{Linux}}
& \multicolumn{2}{c}{\macicon\,\textbf{macOS}}
& \multicolumn{2}{c}{\iosicon\,\textbf{iOS}}
& \multicolumn{2}{c}{\androidicon\,\textbf{Android}}
& \multicolumn{2}{c}{\chromeicon\,\textbf{Web}}
& \multirow{2}{*}{\textbf{Avg.}} \\
\cmidrule(lr){2-3}
\cmidrule(lr){4-5}
\cmidrule(lr){6-7}
\cmidrule(lr){8-9}
\cmidrule(lr){10-11}
\cmidrule(lr){12-13}
& \textbf{Easy} & \textbf{Hard}
& \textbf{Easy} & \textbf{Hard}
& \textbf{Easy} & \textbf{Hard}
& \textbf{Easy} & \textbf{Hard}
& \textbf{Easy} & \textbf{Hard}
& \textbf{Easy} & \textbf{Hard}
& \\
\midrule

\textbf{Human Performance}
& 100.0 & 78.4
& 100.0 & 74.6
& 100.0 & 79.8
& 100.0 & 75.2
& 100.0 & 72.6
& 100.0 & 76.9
& \cellcolor[HTML]{DFF3E4}88.1 \\

\midrule
\rowcolor{modelgroup}
\multicolumn{14}{c}{\textbf{\textit{General Models}}} \\
\midrule

\osmodelicon{openai.pdf}GPT-5.5
& \textbf{84.2} & \textbf{23.7}
& \underline{65.0} & 17.5
& \textbf{81.3} & \textbf{15.6}
& \underline{75.0} & 8.3
& \textbf{83.3} & \textbf{16.7}
& \textbf{78.1} & \underline{16.5}
& \cellcolor[HTML]{E7F6EA}\textbf{47.1} \\

\osmodelicon{openai.pdf}GPT-5.4
& 26.3 & \underline{13.2}
& \textbf{70.0} & \underline{22.5}
& 55.7 & \underline{10.8}
& \textbf{83.3} & \textbf{29.2}
& 33.3 & 0.0
& \underline{55.2} & 15.8
& \cellcolor[HTML]{F0F7E6}\underline{34.6} \\

\osmodelicon{claude-logo.pdf}Claude Opus 4.7
& \underline{47.4} & 5.3
& \underline{65.0} & \textbf{32.5}
& \underline{62.5} & 9.4
& 50.0 & \underline{16.7}
& \underline{41.7} & \underline{12.5}
& 54.9 & \textbf{17.2}
& \cellcolor[HTML]{F0F7E6}\underline{34.6} \\

% Claude Opus 4.6
% & -- & --
% & -- & --
% & -- & --
% & -- & --
% & -- & --
% & -- & --
% & -- \\

\osmodelicon{gemini.pdf}Gemini 3.1 Pro
& 29.2 & 10.8
& 30.0 & 12.5
& 30.4 & 9.6
& 50.0 & 14.1
& \underline{41.7} & 8.4
& 31.6 & 12.3
& \cellcolor[HTML]{FFF6D8}23.4 \\

\osmodelicon{gemini.pdf}Gemini 3 Pro
& 24.9 & 8.6
& 31.2 & 11.3
& 25.1 & 7.4
& 33.3 & 12.5 
& 25.0 & 4.2
& 30.8 & 11.3
& \cellcolor[HTML]{FDEAD2}18.8 \\

\osmodelicon{qwen.pdf}Qwen 3.6 Plus
& 20.5 & 5.1
& 20.0 & 0.0
& 21.6 & 4.3
& 33.3 & 12.5
& 16.7 & 4.2
& 26.6 & 9.2
& \cellcolor[HTML]{F9DCD5}14.5 \\

\osmodelicon{kimi.png}Kimi K2.6
& 5.3 & 0.0
& 15.0 & 2.5
& 6.3 & 0.0
& 8.1 & 0.7
& 6.2 & 0.5
& 12.4 & 1.8
& \cellcolor[HTML]{F2B8B5}4.9 \\

\midrule
\rowcolor{modelgroup}
\multicolumn{14}{c}{\textbf{\textit{Specialized CUA Models}}} \\
\midrule

\osmodelicon{sh_ai_logo.pdf}ScaleCUA-32B
& 15.8 & 0.0
& 40.0 & 2.5
& 0.0 & 0.0
& 8.3 & 0.0
& 0.0 & 0.0
& 25.9 & 4.2
& \cellcolor[HTML]{F5C9C5}8.1 \\

\osmodelicon{hku.pdf}OpenCUA-72B
& 15.8 & 0.0
& 35.0 & 5.0
& 37.5 & 0.0
& 8.3 & 0.0
& 0.0 & 0.0
& 18.1 & 1.5
& \cellcolor[HTML]{F7D1CC}10.1 \\

\bottomrule
\end{tabular}
}
\end{table*}


\subsection{Main Results}

\noindent
\textbf{Current agents remain far from generalist CUA.}
Table~\ref{tab:main_results} shows that even the strongest current models remain far from reliable generalist computer-using agents. GPT-5.5 obtains the best overall average success rate, but still reaches only 47.1\% across platform--difficulty settings. The gap is especially clear on hard tasks, where success rates frequently fall below 20\% and sometimes approach zero. These results indicate that current agents can solve a portion of visually explicit and short-horizon tasks, but still lack the robustness needed for long-horizon planning, precise GUI grounding, state tracking, and final-state persistence in real operating systems.

\noindent
\textbf{Humans are consistent across platforms, while agents are not.}
The human reference row is consistent across all platforms, whereas agent performance varies sharply by operating system and remains far below the human level. For example, GPT-5.4 performs strongly on Linux and iOS tasks, but its Windows and Android tasks success are only 26.3\% and 33.3\%, respectively. Similar drops appear for other models, suggesting that the main challenge is not only performing well on specific platform, but also consistently success on every platforms.

% \noindent
% \textbf{Different models exhibit different platform strengths.}
% The results also reveal model-specific platform strengths rather than a single universal ranking across all settings. GPT-5.5 is the strongest overall and is particularly effective on easy Windows, macOS, Android, and Web tasks. GPT-5.4 is comparatively strong on Linux and iOS, achieving 70.0\% and 83.3\% on easy tasks and the best iOS hard-task result at 29.2\%. Claude Opus 4.7 is strongest on Linux hard tasks, reaching 32.5\%, but is weaker on Windows hard tasks. Gemini 3.1 Pro shows better mobile easy-task performance than its desktop performance, while Qwen 3.6 Plus also performs relatively better on iOS than on desktop and Android. Specialized CUA models do not uniformly outperform general-purpose models: ScaleCUA averages 8.1\% and OpenCUA averages 10.1\%, with isolated strengths on Linux and macOS but weak hard-task performance across most platforms. These trends suggest that progress toward generalist CUA requires not just stronger base models, but also better adaptation to heterogeneous platform conventions, input modalities, and application ecosystems.


\subsection{Analysis}

\noindent
\textbf{Cross-application workflows are harder than single-application commands.}
Cross-application tasks add a second layer of difficulty: the agent must preserve information across window switches, identify which application owns each subgoal, and ensure that all destination states are updated before stopping. As shown in Appendix~\ref{fig:linux_error_case_1}, GPT-5.5 trajectory audit shows that cross-application tasks require substantially longer interaction traces than ordinary single-application tasks on average. The failures are also qualitatively different: agents often complete a visually plausible intermediate action, but fail the final evaluator because a calendar event or follow-up message is missing or inconsistent. This makes cross-application evaluation closer to real-world computer use, where user intent is rarely confined to a single GUI command.

\noindent
\textbf{Cross-platform settings introduce additional transfer difficulty.}
Cross-platform tasks further stress agents because they require correct use of platform-specific interaction conventions and reliable transfer of state or artifacts between environments. In single-platform tasks, the agent can remain within one action backend and one set of visual conventions. Cross-platform workflows instead require the agent to create a source-side artifact, hand it off through the runtime, and consume it correctly in a target environment.  As shown in Appendix~\ref{fig:linux_error_case_1}, the inspected GPT-5.5 trajectories show this extra difficulty clearly. Together with the platform sensitivity observed in Table~\ref{tab:main_results}, this suggests that cross-platform CUA is not simply a larger task set, but a distinct generalization problem involving heterogeneous UI conventions and input primitives.

\noindent
\textbf{Error analysis.}
We further inspect trajectories to categorize common failure modes. The most frequent errors fall into four groups. First, \emph{visual grounding errors} occur when the agent clicks a nearby control or misreads dense document and diagram interfaces. Second, \emph{state-tracking errors} appear when the agent loses copied information across applications or forgets a required destination file. Third, \emph{persistence errors} occur when the visible interface looks correct but the durable final state is not saved or committed in the format expected by the evaluator. Fourth, \emph{platform and timing errors} arise from modal dialogs, delayed application loading or missing confirmation steps. These categories explain why screenshot-only agents can often produce plausible trajectories while still failing the task. A complete qualitative error study is provided in Appendix~\ref{appendix:common_errors}.
